\chapter*{ЗАКЛЮЧЕНИЕ}
\addcontentsline{toc}{chapter}{ЗАКЛЮЧЕНИЕ}

3Д--реконструкция, благодаря возможностям точного восстановления объектов и сцен в трёхмерном пространстве,
является востребованным направлением и активно используется в различных областях,
таких как робототехника, медицина, игры, киноиндустрия и многие другие.
Интерес к этой темы выявил, что распространена лишь одна классификация методов 3Д-восстановления:
разделение на активные и пассивные подходы~\cite{stereosop2, sfm, structLight, ref1}.
С целью создание классификации, охватывающей большее количество характеристик и
была подготовлена данная научно-исследовательская работа, в ходе которой:

\begin{itemize}
    \item проведен анализ предметной области «обработка изображений» с акцентом на конвертации плоских снимков в 3Д модели;
    \item описаны основные методы конвертации 2Д изображений в 3Д, а именно: метод структурного освещения, глубины из размытия, стереосопоставления, структуры из движения и комфортной генерации глубины. Каждому методу приведены основные шаги их использования, особенности, и ограничения области применения;
    \item выделены критерии сравнения методов, и необходимые для реализации разных типов проектов;
    \item создана классификация описанных методов по выделенным критериям, которая необходима при выборе наиболее подходящего метода для конкретной задачи и значительно упрощает процесс принятия решений.
\end{itemize}
Таким образом, можно констатировать достижение поставленной цели научно-исследовательской работы.

 

 

